\documentclass[journal]{IEEEtran}
\usepackage{amsmath,amsfonts,amssymb,amsthm}
\usepackage{algorithmic}
\usepackage{algorithm}
\usepackage{array}
\usepackage{textcomp}
\usepackage{stfloats}
\usepackage{url}
\usepackage{verbatim}
\usepackage{graphicx}
\usepackage{cite}
\usepackage{booktabs}
\usepackage{multirow}
\usepackage{hyperref}

\hypersetup{
    colorlinks=true,
    linkcolor=black,
    citecolor=black,
    urlcolor=blue
}

\newtheorem{theorem}{Theorem}
\newtheorem{lemma}{Lemma}
\newtheorem{definition}{Definition}

\begin{document}

\title{Adaptive Lookahead Routing, Fault-Tolerant Live Detours, and Belief-State Filtering in Energy-Harvesting Wireless Sensor Networks}

\author{Ananthakrishnan~R,~\IEEEmembership{Student~Member,~IEEE,}
        Santhosh~G,~\IEEEmembership{Student~Member,~IEEE,}
        and~Thilang~TV,~\IEEEmembership{Student~Member,~IEEE}%
\thanks{The authors are with the Department of Electrical and Computer Science Engineering, School of Electrical Engineering (SELECT), Vellore Institute of Technology (VIT), Vellore 632014, Tamil Nadu, India (e-mail: ananthakrishnan.r2025@vitstudent.ac.in; santhosh.g2025@vitstudent.ac.in; thilang.tv2025@vitstudent.ac.in).}%
\thanks{Author Program \& Registration Numbers: B.Tech in Electrical and Computer Science Engineering --- Ananthakrishnan R (25BEL0005), Santhosh G (25BEL0016), Thilang TV (25BEL0036).}%
\thanks{Open-source code, testbenches, and interactive simulation dashboard are publicly accessible at GitHub: \url{https://github.com/santhoshvellore7119-web/WSN_proj}.}%
}

\markboth{IEEE Transactions on Mobile Computing / IEEE Sensors Journal, Draft Paper 2026}%
{Ananthakrishnan \MakeLowercase{\textit{et al.}}: Adaptive Lookahead Routing \& Belief-State Filtering in EH-WSNs}

\maketitle

\begin{abstract}
Wireless Sensor Networks (WSNs) powered by ambient energy harvesting (solar irradiance, thermal gradients, and ambient RF) violate the classical assumption of monotonic electrochemical battery depletion ($E_i(t_1) \ge E_i(t_2)$ for $t_1 \le t_2$). Standard routing protocols---including Dijkstra's shortest path, standard LEACH dynamic clustering, and static bottleneck Bellman-Ford---evaluate candidate relay nodes based on static energy snapshots captured at round inception ($t=0$). This induces two severe structural pathologies: \textit{false rejection} of recharging relays (nodes whose initial low residual energy is rapidly replenished by ambient influx during multi-hop packet transit delay) and \textit{accelerated burnout} of occluded, non-harvesting nodes. Furthermore, when continuous battery telemetry is suppressed to conserve radio power, hidden mid-round relay brownouts trigger catastrophic in-flight packet drops.

In this paper, we propose a comprehensive, lightweight algorithmic framework for Energy-Harvesting WSNs (EH-WSNs) grounded in fundamental Data Structures and Algorithms (DSA):
(1) \textbf{3D Time-Augmented Dynamic Programming ($dp[v][h][t]$):} A polynomial-time routing algorithm over a time-expanded Directed Acyclic Graph (DAG) that explicitly plans multi-hop routes around projected solar replenishment, provably guaranteeing a $2\epsilon$-approximation bound under bounded stochastic arrival variance $|\xi| \le \epsilon$.
(2) \textbf{Disjoint-Set Union (DSU) Live Detour Recovery:} An in-flight fault-repair algorithm using Union-Find with path compression and rank optimization that splices detour bridges around mid-round brownout relays in $O(\text{deg}(u)\alpha(V))$ time, delivering a $3.8\times - 6.1\times$ execution speedup over full DP recomputation with $0\%$ packet loss up to a $30\%$ node failure rate.
(3) \textbf{Probabilistic Belief Routing Under Partial Observability:} An architectural decoupling of physical ground truth from router observability via a 24-bin discrete Histogram Belief Filter with exact boundary clipping at 0 and $C_i$, coupled with log-domain Max-Product Belief DP with recharge waiting deadlines ($w \in [0, D]$) and sequential load reservations.
(4) \textbf{Harvesting-Aware Dynamic Clustering (EH-LEACH):} Solar-weighted cluster head elections with min-heap priority queue tiebreaking.

We validate the framework through an automated 68-test unit verification suite (100\% pass rate) and multi-seed statistical benchmarks ($N=30$ seeds, 350 rounds). Empirical results confirm that under scarce energy regimes, our belief-driven routing closes up to $81.5\%$ of the omniscient oracle delivery gap while reducing Expected Calibration Error (ECE = 0.019 vs 0.039) and Brier scores over stale telemetry. Furthermore, we characterize the first-order radio electronic reception dissipation penalty ($E_{\text{rx}} = k \cdot E_{\text{elec}}$), delineating exact switching boundaries where direct single-hop transmission outperforms multi-hop relaying.
\end{abstract}

\begin{IEEEkeywords}
Energy Harvesting, Wireless Sensor Networks, Dynamic Programming, Disjoint-Set Union, Fault-Tolerant Routing, Partially Observable Markov Decision Processes, Telemetry Calibration, Radio Dissipation Physics.
\end{IEEEkeywords}

\section{Introduction}

\IEEEPARAGRAPHstart{W}{ireless} Sensor Networks (WSNs) are deployed extensively across environmental telemetry, precision agriculture, structural health monitoring, and perimeter defense \cite{akyildiz2002survey, sah2020renewable}. Historically, sensor motes have been constrained by primary electrochemical batteries (e.g., Lithium coin cells). In battery-powered networks, stored energy decreases monotonically with every sensing, computation, and transmission event:
\begin{equation}
E_i(t_1) \ge E_i(t_2) \quad \forall\ t_1 \le t_2
\end{equation}
Consequently, classical routing protocols---such as Dijkstra's shortest path algorithm, A* search, standard LEACH dynamic clustering \cite{heinzelman2002application}, and static bottleneck Bellman-Ford---operate under the implicit invariant that a node with low residual energy is permanently decaying. To prevent in-flight packet drops, these protocols rigidly prune low-battery nodes from forwarding graphs.

Recent engineering breakthroughs in micro-photovoltaic cells, piezoelectric harvesters, and ambient RF rectennas have enabled modern sensor motes to operate perpetually via ambient energy harvesting \cite{sudevalayam2011energy, sharma2020energy}. This introduces a fundamental paradigm shift: \textbf{node battery capacity is non-monotonic}. A candidate relay node that possesses depleted reserves at the beginning of a transmission cycle can harvest substantial replenishment while preceding data packets traverse intermediate multi-hop hops.

\subsection{Pathologies in Existing State-of-the-Art}
When classical static routing protocols are deployed in renewable energy-harvesting sensor networks, their blindness to incoming energy induces two primary failure modes:

\begin{enumerate}
    \item \textbf{The False Rejection Pathology:} Consider a multi-hop path where source cluster head $s$ routes data to base station $v_{\text{BS}}$ through candidate relay $u$. At round inception ($t = 0$), relay $u$ has residual battery $E_u(0) < E_{\text{tx}}$. Classical snapshot algorithms immediately prune node $u$. However, because packet forwarding across preceding hops incurs non-negligible transmission, queuing, and TDMA processing delays ($t_i = i \cdot \delta$), node $u$ receives ambient solar irradiance during transit. By the time the packet physically arrives at node $u$, its battery has replenished to operational capacity. Discarding node $u$ forces the network onto circuitous detours that burn excess global energy.
    \item \textbf{The Shaded Burnout Pathology:} In real-world agricultural or forest deployments, solar exposure is highly heterogeneous. Nodes situated beneath dense canopy or building occlusions receive zero solar recharge ($H_v(t) \approx 0$), yet may begin with substantial initial electrochemical reserves. Snapshot protocols repeatedly select these shaded nodes as intermediate forwarders due to their high instantaneous battery levels. Consequently, shaded nodes burn out rapidly, severing network connectivity while sunny peripheral nodes sit idle with saturated batteries.
    \item \textbf{In-Flight Relay Brownout \& Inefficient Recovery:} Environmental energy arrivals are inherently stochastic. When an intermediate relay node experiences a sudden cloud occlusion or voltage drop during packet transit, in-flight data packets are dropped. Existing protocols either accept high packet loss or invoke global routing recalculations across the entire network ($O(|E|HT)$), introducing prohibitive latency and control overhead.
\end{enumerate}

\subsection{Major Contributions}
To overcome these challenges without incurring the high memory and computational overhead of Deep Reinforcement Learning (DRL) or integer linear programming, this paper develops a deterministic, polynomial-time framework grounded in fundamental Data Structures and Algorithms:

\begin{enumerate}
    \item \textbf{3D Time-Augmented Dynamic Programming ($dp[v][h][t]$):} We formulate a backward dynamic programming recurrence over a time-expanded DAG $\mathcal{G}_T$ that plans routes across node, hop count, and time dimensions. We formally prove that under bounded stochastic arrival variance $|\xi| \le \epsilon$, our algorithm guarantees a $2\epsilon$-approximation bound relative to an omniscient offline oracle.
    \item \textbf{Disjoint-Set Union (DSU) Live Detour Recovery:} We design an in-flight fault recovery algorithm utilizing Disjoint-Set Union with path compression and union-by-rank. By maintaining connected components across operable nodes, an upstream predecessor identifies an alternate bridge node in $O(\text{deg}(u)\alpha(V))$ time, providing a $3.8\times - 6.1\times$ latency reduction over global DP recomputation with zero packet loss.
    \item \textbf{Probabilistic Belief Routing Under Partial Observability:} We formulate an architectural truth/belief split via an \texttt{EnergyView} protocol. Router beliefs are tracked through a 24-bin discrete Histogram Belief Filter with exact boundary clipping at 0 (brownout) and $C_i$ (saturation), eliminating Gaussian truncation distortion. We solve for routes and waiting schedules using log-domain Max-Product Belief DP with sequential load reservations.
    \item \textbf{Relay Reception Penalty Characterization ($E_{\text{rx}}$):} We incorporate first-order electronic receiver dissipation ($E_{\text{rx}} = k \cdot E_{\text{elec}}$), demonstrating both analytically and empirically that multi-hop relaying is not universally energy-saving. We define the exact crossover boundaries where single-hop direct transmission preserves more aggregate network energy.
    \item \textbf{Full-Stack Implementation \& Multi-Seed Benchmarks:} We deploy a complete software suite comprising a Python simulation engine, an asynchronous FastAPI backend, an interactive web dashboard, 68 passing automated unit tests, and multi-seed statistical evaluations ($N=30$ seeds) across solar, shaded, and stochastic regimes.
\end{enumerate}

\section{Related Work \& Literature Review}

Routing in energy-harvesting sensor networks has been investigated across multiple paradigms. Table~\ref{tab:literature} summarizes landmark publications from 2020--2025 across IEEE, ACM, and Elsevier, detailing their methodologies and critical research gaps.

\begin{table*}[!t]
\caption{Comprehensive Prior Art Matrix and Comparison with Proposed Framework (2020--2025)}
\label{tab:literature}
\centering
\begin{tabular}{@{}llllp{5.8cm}@{}}
\toprule
\textbf{Reference \& Authors} & \textbf{Venue \& Year} & \textbf{Paradigm} & \textbf{Complexity} & \textbf{Critical Limitations \& Research Gaps} \\
\midrule
B. Han \textit{et al.} \cite{han2022novel} & \textit{MDPI Sensors} (2022) & Solar EH-LEACH & $O(N)$ & Ignores packet transit delay ($\delta$); cannot model recharge occurring during forwarding. \\
A. Naeem \textit{et al.} \cite{naeem2021dare} & \textit{IEEE Trans. Green Commun.} (2021) & DARE-SEP Clustering & $O(N \log N)$ & Assumes non-rechargeable batteries; lacks live fault recovery under mid-round relay depletion. \\
D. K. Sah \& T. Amgoth \cite{sah2023harvested} & \textit{IEEE Sensors J.} (2023) & EWMA Solar Dijkstra & $O(|E| + |V|\log |V|)$ & Single-path optimization; prone to false rejection of transient recharging relays. \\
R. Tan \textit{et al.} \cite{tan2021throughput} & \textit{ACM Trans. Sen. Netw.} (2021) & Time-expanded ILP & NP-hard & Centralized offline batch solver; cannot execute on distributed 8/32-bit motes. \\
J. Kim \& S. Park \cite{kim2021energy} & \textit{IEEE Internet of Things J.} (2021) & Opportunistic ETX & $O(|E|)$ & Triggers expensive global table recalculation upon relay failure; lacks local detours. \\
M. A. Aref \& S. Jayaweera \cite{aref2023deep} & \textit{IEEE Trans. Cogn. Commun.} (2023) & Deep Q-Networks (DQN) & Matrix mult. & Memory footprint exceeds 1 MB; computationally infeasible for low-power MCUs. \\
M. Rathee \textit{et al.} \cite{rathee2021qos} & \textit{IEEE Trans. Eng. Manage.} (2021) & Ant Colony Opt. (ACO) & Non-deterministic & Meta-heuristic; lacks formal approximation bounds under bounded harvest variance. \\
Y. Zhang \textit{et al.} \cite{zhang2021energy} & \textit{IEEE Access} (2021) & Solar Cluster Radius & $O(N)$ & Assumes single-hop intra-cluster delivery; ignores receiver dissipation ($E_{\text{rx}}$). \\
S. K. Sharma \textit{et al.} \cite{sharma2020energy} & \textit{IEEE Commun. Surv. Tutorials} (2020) & Comprehensive Survey & N/A & Identifies transit lookahead routing and live fault tolerance as critical open gaps. \\
\midrule
\textbf{Proposed Framework} & \textbf{This Work (2026)} & \textbf{3D Time-DP + DSU} & \textbf{$O(|E|HT)$} & \textbf{Recharge lookahead, $O(\text{deg}(u)\alpha(V))$ live DSU detours, $<50$ KB RAM.} \\
\bottomrule
\end{tabular}
\end{table*}

\begin{table}[!t]
\caption{Methodology Comparison Across Algorithmic Dimensions}
\label{tab:method_comparison}
\centering
\begin{tabular}{@{}lccc@{}}
\toprule
\textbf{Dimension} & \textbf{Classical} & \textbf{DRL / Swarm} & \textbf{Proposed (Ours)} \\
\midrule
Recharge Lookahead & None (Static $t=0$) & Implicit & \textbf{Exact Table ($dp[v][h][t]$)} \\
Algorithmic Time & $O(|E| + |V|\log |V|)$ & Non-poly & \textbf{$O(|E| \cdot H \cdot T)$} \\
RAM Footprint & $< 5$ KB & $> 1$ MB & \textbf{$< 50$ KB (MCU Viable)} \\
Fault Recovery & Drop / Recompute & Retrain & \textbf{$O(\text{deg}(u)\alpha(V))$ DSU} \\
Theoretical Bounds & Dijkstra static & None & \textbf{Provable $2\epsilon$-bound} \\
\bottomrule
\end{tabular}
\end{table}

\section{System Model \& Problem Formulation}

\subsection{Network Architecture}
We model a wireless sensor network deployed over a 2D field of dimensions $W \times H_{\text{field}}$ as a directed communication graph $G = (V, E)$:
\begin{itemize}
    \item $V = \{v_1, v_2, \dots, v_N\}$ is the set of $N$ stationary sensor motes.
    \item $v_{\text{BS}} \equiv -1$ denotes the central Base Station (sink) located at coordinates $(x_{\text{BS}}, y_{\text{BS}})$.
    \item A directed edge $(u, v) \in E$ exists if the Euclidean distance $d(u, v) = \sqrt{(x_u - x_v)^2 + (y_u - y_v)^2} \le R_{\text{tx}}$, where $R_{\text{tx}}$ is the maximum radio transmission range.
\end{itemize}

Each sensor node $v \in V$ is characterized by battery capacity $C_v = E_{\text{max}}$, instantaneous residual battery $E_v(t) \in [0, E_{\text{max}}]$, ambient energy harvesting rate $H_v(t) \ge 0$, battery leakage rate $\lambda$, and charging efficiency $\eta \in (0, 1]$. Continuous battery evolution satisfies:
\begin{equation}
\frac{dE_v(t)}{dt} = \min\left( \eta H_v(t) - \lambda E_v(t) - P_{\text{drain}}(v, t), \, \frac{E_{\text{max}} - E_v(t)}{\Delta t} \right)
\end{equation}
subject to the brownout boundary condition $E_v(t) \ge 0$.

\subsection{First-Order Radio Dissipation Physics}
Energy dissipation per transmission and reception adheres to Heinzelman's first-order radio model \cite{heinzelman2002application}:
\begin{equation}
E_{\text{tx}}(k, d) = \begin{cases} 
k \cdot E_{\text{elec}} + k \cdot \epsilon_{\text{fs}} \cdot d^2, & d < d_0 \quad (\text{Free-space}) \\ 
k \cdot E_{\text{elec}} + k \cdot \epsilon_{\text{mp}} \cdot d^4, & d \ge d_0 \quad (\text{Multipath}) 
\end{cases}
\end{equation}
\begin{equation}
E_{\text{rx}}(k) = k \cdot E_{\text{elec}}
\end{equation}
The crossover distance threshold $d_0$ is defined as:
\begin{equation}
d_0 = \sqrt{\frac{\epsilon_{\text{fs}}}{\epsilon_{\text{mp}}}} = \sqrt{\frac{10 \times 10^{-12}\text{ J/bit/m}^2}{0.0013 \times 10^{-12}\text{ J/bit/m}^4}} \approx 87.70\text{ meters}
\end{equation}
where hardware parameters calibrated from the TI CC2420 transceiver are: $E_{\text{elec}} = 50\text{ nJ/bit}$, $\epsilon_{\text{fs}} = 10\text{ pJ/bit/m}^2$, $\epsilon_{\text{mp}} = 0.0013\text{ pJ/bit/m}^4$, and packet size $k = 4000\text{ bits}$ (500 bytes).

\begin{table}[!t]
\caption{System Simulation & Transceiver Parameters}
\label{tab:params}
\centering
\begin{tabular}{@{}lll@{}}
\toprule
\textbf{Parameter} & \textbf{Symbol} & \textbf{Physical Value} \\
\midrule
Network Field Size & $W \times H_{\text{field}}$ & $100\text{ m} \times 100\text{ m}$ \\
Total Sensor Motes & $N$ & $50$ motes \\
Base Station Sink & $v_{\text{BS}}$ & $(50.0, 50.0)$ \\
Radio Transmission Range & $R_{\text{tx}}$ & $48.0\text{ m}$ \\
Electronic Circuit Dissipation & $E_{\text{elec}}$ & $50\text{ nJ/bit}$ \\
Free-Space Amplifier Const. & $\epsilon_{\text{fs}}$ & $10\text{ pJ/bit/m}^2$ \\
Multipath Amplifier Const. & $\epsilon_{\text{mp}}$ & $0.0013\text{ pJ/bit/m}^4$ \\
Crossover Distance Threshold & $d_0$ & $87.70\text{ m}$ \\
Data Packet Payload & $k$ & $4000\text{ bits}$ (500 B) \\
Initial Battery Capacity & $C_{\text{max}}$ & $0.50\text{ J}$ \\
Receiver Penalty per Hop & $E_{\text{rx}}$ & $0.20\text{ mJ}$ \\
\bottomrule
\end{tabular}
\end{table}

\subsection{The Electronic Reception Overhead ($E_{\text{rx}}$) Trade-Off}
In multi-hop routing, each intermediate forwarder incurs electronic reception dissipation $E_{\text{rx}} = k \cdot E_{\text{elec}} = 0.20\text{ mJ}$. Consider node $u$ communicating with base station $v_{\text{BS}}$ over distance $d = 60\text{ m}$ ($d < d_0$):
\begin{itemize}
    \item \textbf{Single-hop Direct:} $E_{\text{direct}} = 4000 \cdot (50\text{ nJ} + 10\text{ pJ} \cdot 60^2) = 0.344\text{ mJ}$.
    \item \textbf{Two-hop Multi-hop ($d_1 = 30\text{m}, d_2 = 30\text{m}$):} Transmit costs are $E_{\text{tx1}} = 0.236\text{ mJ}$, $E_{\text{tx2}} = 0.236\text{ mJ}$, plus intermediate reception $E_{\text{rx}} = 0.200\text{ mJ}$. Total network dissipation $= 0.672\text{ mJ}$.
\end{itemize}
Thus, multi-hop relaying is not unconditionally optimal; it is advantageous only when the transmission distance reduction ($d^4$ vs $d^2$) or intermediate ambient solar harvesting overcomes the fixed electronic reception overhead $E_{\text{rx}}$.

\subsection{Optimization Formulation}
Let $\mathcal{P}(s, v_{\text{BS}})$ be the set of loop-free paths from source cluster head $s$ to sink $v_{\text{BS}}$. With inter-hop delay $\delta$, arrival epoch at hop $i$ is $t_i = t_{\text{curr}} + i \cdot \delta$. The min-max bottleneck lifetime objective is:
\begin{equation}
P^* = \arg\max_{P \in \mathcal{P}(s, v_{\text{BS}})} \min_{0 \le i < |P|} E_{\text{projected}}(v_i, t_{\text{curr}} + i \cdot \delta)
\end{equation}
subject to: (1) $|P| \le H$, (2) $E_{\text{projected}}(v_i, t_i) \ge E_{\text{tx}}(k, d(v_i, v_{i+1}))$, and (3) $v_i \ne v_j\ \forall\ i \ne j$.

\section{Proposed Algorithmic Framework}

\subsection{3D Time-Augmented Dynamic Programming}
We define the 3D dynamic programming table $dp[v][h][t]$ representing the maximum bottleneck residual energy along a path reaching node $v$ after $h$ hops at time offset $t$:
\begin{equation}
dp[v][h][t] = \max_{u \in \text{nbr}(v)} \min\left( dp[u][h-1][t-\delta], \, E_{\text{proj}}(v, t_{\text{curr}} + t) - E_{\text{rx}} \right)
\end{equation}
subject to $dp[u][h-1][t-\delta] \ge E_{\text{tx}}(k, d(u, v))$, with projected battery bounded by capacity:
\begin{equation}
E_{\text{proj}}(v, t_{\text{curr}} + t) = \min\left( C_v, \, E_v(t_{\text{curr}}) + \int_{t_{\text{curr}}}^{t_{\text{curr}} + t} H_v(\tau) d\tau \right)
\end{equation}
Base initialization: $dp[s][0][0] = E_s(t_{\text{curr}})$, and $dp[v][h][t] = -\infty$ elsewhere. Base station sink $v_{\text{BS}} = -1$ has $C_{\text{BS}} = \infty$ and $E_{\text{rx}} = 0$.

\begin{theorem}[Approximation Bound Under Bounded Variance]
Let $P^*$ be the optimal path selected by an omniscient oracle with complete future knowledge of stochastic arrivals. Under bounded stochastic harvest error $|\xi_v(t)| \le \epsilon$, the path $P_{\text{Time-DP}}$ computed by 3D Time-Augmented DP satisfies:
\begin{equation}
\text{Bottleneck}(P_{\text{Time-DP}}) \ge \text{Bottleneck}(P^*) - 2\epsilon
\end{equation}
\end{theorem}

\begin{proof}
Let $\hat{B}(P)$ denote projected bottleneck energy and $B(P)$ realized physical energy. Since $|\xi_v(t)| \le \epsilon$, for all paths $P$:
\begin{equation}
|B(P) - \hat{B}(P)| \le \epsilon
\end{equation}
Because Time-Augmented DP guarantees exact optimality over projected values:
\begin{equation}
\hat{B}(P_{\text{Time-DP}}) \ge \hat{B}(P^*)
\end{equation}
Applying the triangle inequality:
\begin{align}
B(P_{\text{Time-DP}}) &\ge \hat{B}(P_{\text{Time-DP}}) - \epsilon \\
&\ge \hat{B}(P^*) - \epsilon \\
&\ge (B(P^*) - \epsilon) - \epsilon = B(P^*) - 2\epsilon
\end{align}
This completes the proof.
\end{proof}

\subsection{Disjoint-Set Union (DSU) Live Detour Recovery}
When an intermediate relay node depletes unexpectedly mid-round:
\begin{enumerate}
    \item Upstream predecessor $u = v_{\text{fail}-1}$ detects that next-hop relay $v$ is inoperable.
    \item Predecessor $u$ queries its 1-hop physical neighbors $w \in \text{nbr}(u)$.
    \item Connectivity to base station is validated in $O(\alpha(V))$ via Disjoint-Set Union with path compression and rank optimization:
    \begin{equation}
    \text{find}(w) == \text{find}(v_{\text{BS}})
    \end{equation}
    \item Candidate bridge node $b$ minimizing link cost $E_{\text{tx}}(k, d(u, b))$ is spliced into the active path without resetting global tables.
\end{enumerate}

Splicing candidate verification requires $O(\text{deg}(u)\alpha(V))$ time, achieving a $3.8\times - 6.1\times$ execution speedup over full DP recomputation ($0.92\text{ ms}$ vs $5.13\text{ ms}$) while maintaining $100\%$ packet delivery ratio.

\subsection{Probabilistic Belief-State Filtering}
To prevent heavy telemetry overhead, motes report battery levels every $R$ rounds. The router tracks battery uncertainty via a 24-bin discrete Histogram Belief Filter:
\begin{itemize}
    \item \textbf{Boundary Clipping:} Finite battery bounds at 0 (brownout) and $C_i$ (saturation) accumulate boundary mass, eliminating Gaussian tail truncation distortion.
    \item \textbf{Sheppard's Correction:} Adjusts variance by $\delta_{\text{bin}}^2 / 12$ across discrete state convolutions.
    \item \textbf{Max-Product Log-DP:} Routes maximize log-confidence over path hops and waiting delays $w \in [0, D]$:
    \begin{equation}
    \max_{\pi, w} \sum_{(u, v) \in \pi} \log P(\text{afford hop } (u, v) \text{ at slot } t)
    \end{equation}
\end{itemize}

\subsection{Harvesting-Aware Dynamic Clustering (EH-LEACH)}
Cluster head election thresholds are dynamically scaled by projected solar intake:
\begin{equation}
T_{\text{EH-LEACH}}(n) = T(n) \cdot \left[ 0.5 \cdot \frac{E_{\text{residual}}(n)}{E_{\text{max}}} + 0.5 \cdot \frac{E_{\text{solar\_intake}}(n)}{E_{\text{solar\_peak}}} \right]
\end{equation}
Ties are resolved using a min-heap priority queue in $O(\log K)$ time.

\section{System Architecture \& Software Implementation}

The framework is implemented as an end-to-end, modular software package:
\begin{itemize}
    \item \textbf{Core Simulation Engine (\texttt{src/}):} Implements 1st-order radio dissipation, diurnal solar generation, 3D Time-DP, DSU live detours, and 24-bin histogram belief views.
    \item \textbf{Asynchronous FastAPI Backend (\texttt{backend/}):} REST API exposing \texttt{/simulate}, \texttt{/benchmark}, and \texttt{/experiments/scalability} with SQLite run persistence and threaded task execution.
    \item \textbf{Automated Test Suite (\texttt{tests/}):} 68 automated unit and integration tests verifying all radio models, DSU detours, DP bounds, and API contracts (100\% passing in 5.44 seconds).
\end{itemize}

\begin{table}[!t]
\caption{Automated Pytest Verification Suite Breakdown}
\label{tab:pytest}
\centering
\begin{tabular}{@{}llcc@{}}
\toprule
\textbf{Test Module} & \textbf{Algorithm / Invariant} & \textbf{Tests} & \textbf{Status} \\
\midrule
\texttt{test\_energy\_model.py} & Radio $d^2/d^4$ dissipation, $d_0 \approx 87.7\text{m}$ & 4 & \textbf{PASS} \\
\texttt{test\_adversarial\_counterexample.py} & Recharging path selection & 1 & \textbf{PASS} \\
\texttt{test\_dp\_lifetime.py} & 3D DP table filling \& clamping & 5 & \textbf{PASS} \\
\texttt{test\_rip\_up\_reroute.py} & DSU path compression \& detours & 4 & \textbf{PASS} \\
\texttt{test\_dsu\_speedup.py} & DSU $5.5\times$ latency reduction & 1 & \textbf{PASS} \\
\texttt{test\_uncertainty\_routing.py} & Histogram calibration \& ECE & 7 & \textbf{PASS} \\
\texttt{test\_energy\_view.py} & Oracle, Stale, Gaussian views & 5 & \textbf{PASS} \\
\texttt{test\_clustering.py} & LEACH / EH-LEACH probability & 4 & \textbf{PASS} \\
\texttt{test\_harvesting\_model.py} & Diurnal solar \& canopy shade & 6 & \textbf{PASS} \\
\texttt{test\_backend\_api.py} & FastAPI endpoints \& SQLite & 10 & \textbf{PASS} \\
\texttt{test\_network.py}, Baselines & Graph models \& trace replay & 21 & \textbf{PASS} \\
\midrule
\textbf{Total Testbench} & \textbf{Entire Algorithmic Core} & \textbf{68} & \textbf{100\%} \\
\bottomrule
\end{tabular}
\end{table}

\section{Experimental Results \& Performance Evaluation}

\subsection{Adversarial Counterexample Verification}
On a deterministic 5-node topology (\texttt{run\_counterexample.py}) where Relay A (node 1) starts with moderate energy ($0.030\text{ J}$) but zero recharge ($H_1 = 0$), and Relay B (node 2) starts depleted ($0.005\text{ J}$) with ambient recharge ($H_2 = +0.035\text{ J/step}$):
\begin{itemize}
    \item Classical Dijkstra and static maximin evaluate nodes at $t=0$, greedily choosing Relay A, which causes premature route failure at step 8.
    \item 3D Time-Augmented DP correctly projects that Relay B replenishes to $0.0400\text{ J}$ by arrival epoch $t=1$, safely routing along $0 \to 2 \to \text{BS}$ with projected bottleneck residual energy of $0.0400\text{ J}$ and arrival schedule $[0, 1, 2]$, achieving 100\% packet delivery.
\end{itemize}

\begin{table*}[!t]
\caption{Multi-Seed Statistical Benchmark Results Across Regimes ($N = 30$ Seeds, 350 Rounds, Mean $\pm$ Std [95\% CI])}
\label{tab:multi_seed}
\centering
\begin{tabular}{@{}lccc@{}}
\toprule
\textbf{Simulation Regime \& Configuration} & \textbf{First Node Death (FND)} & \textbf{Alive Nodes (Round 350)} & \textbf{Total Residual Energy (Joules)} \\
\midrule
1. Baseline (No Harvesting) & $89.2 \pm 5.4$ & $0.4 \pm 0.5$ & $0.0001 \pm 0.0002$ $[0.0000, 0.0002]$ \\
2. Solar (Unaware LEACH + Dijkstra) & $351.0 \pm 0.0$ & $50.0 \pm 0.0$ & $\mathbf{1.7050 \pm 0.0813}$ $[1.6759, 1.7341]$ \\
3. Solar (Adaptive Time-DP) & $351.0 \pm 0.0$ & $50.0 \pm 0.0$ & $1.6299 \pm 0.0939$ $[1.5963, 1.6635]$ \\
4. Shadowed Solar (Unaware Dijkstra) & $\mathbf{110.2 \pm 9.6}$ & $\mathbf{24.1 \pm 3.2}$ & $\mathbf{0.3900 \pm 0.1996}$ $[0.3186, 0.4614]$ \\
5. Shadowed Solar (Time-DP) & $109.0 \pm 9.5$ & $23.7 \pm 3.8$ & $0.3521 \pm 0.1958$ $[0.2820, 0.4222]$ \\
6. Stochastic Poisson (Unaware LEACH) & $351.0 \pm 0.0$ & $50.0 \pm 0.0$ & $\mathbf{15.8599 \pm 0.1470}$ $[15.8073, 15.9125]$ \\
7. Stochastic Poisson (Adaptive Time-DP) & $351.0 \pm 0.0$ & $50.0 \pm 0.0$ & $15.7890 \pm 0.1527$ $[15.7343, 15.8437]$ \\
\bottomrule
\end{tabular}
\end{table*}

\begin{table*}[!t]
\caption{Energy Uncertainty Routing Benchmark Across Solar Regimes (180 Multi-Seed Runs)}
\label{tab:uncertainty}
\centering
\begin{tabular}{@{}llcccccc@{}}
\toprule
\textbf{Solar Regime} & \textbf{Routing Method} & \textbf{Packet Delivery Ratio} & \textbf{Hop Latency} & \textbf{Wasted Energy} & \textbf{Brier Score} & \textbf{ECE} & \textbf{Oracle Gap Closed} \\
\midrule
\multirow{5}{*}{\textbf{Poor Regime}} 
 & M0 (Blind Dijkstra) & $47.50\% \pm 2.77\%$ & $2.62$ & $0.0268\text{ J}$ & $0.0262$ & $0.0262$ & Baseline \\
 & M1 (Stale Dijkstra, $R=20$) & $46.03\% \pm 3.76\%$ & $2.40$ & $0.0212\text{ J}$ & $0.0394$ & $0.0394$ & $-17.3\%$ \\
 & \textbf{M3 (Belief DP, Proposed)} & $\mathbf{47.47\% \pm 2.29\%}$ & $2.47$ & $0.0227\text{ J}$ & $0.0272$ & $0.0230$ & $\mathbf{+1.44\% \text{ vs M1}}$ \\
 & \textbf{M4 (Belief LCB-Maximin)} & $\mathbf{52.19\% \pm 2.19\%}$ & $3.04$ & $0.0216\text{ J}$ & $0.0233$ & $0.0192$ & $\mathbf{+55.4\% \text{ Gap Closed}}$ \\
 & M2 (Oracle Upper Bound) & $55.96\% \pm 2.72\%$ & $3.46$ & $0.0145\text{ J}$ & $0.0000$ & $0.0000$ & $100.0\%$ \\
\midrule
\multirow{4}{*}{\textbf{Balanced Regime}}
 & M0 (Blind Dijkstra) & $60.73\% \pm 3.24\%$ & $2.66$ & $0.0088\text{ J}$ & $0.0049$ & $0.0049$ & Baseline \\
 & M1 (Stale Dijkstra, $R=10$) & $67.81\% \pm 2.44\%$ & $2.69$ & $0.0009\text{ J}$ & $0.0046$ & $0.0046$ & $+120.1\%$ \\
 & \textbf{M3 (Belief DP, Proposed)} & $\mathbf{65.53\% \pm 3.42\%}$ & $2.70$ & $0.0027\text{ J}$ & $0.0019$ & $0.0022$ & $\mathbf{+81.5\% \text{ Gap Closed}}$ \\
 & M2 (Oracle Upper Bound) & $66.63\% \pm 3.20\%$ & $3.33$ & $0.0000\text{ J}$ & $0.0000$ & $0.0000$ & $100.0\%$ \\
\bottomrule
\end{tabular}
\end{table*}

\subsection{Multi-Seed Benchmark Results}
Table~\ref{tab:multi_seed} compiles 30-seed evaluations across 350 rounds. In unconstrained solar regimes, network lifetime persists indefinitely (50 alive nodes at round 350). In canopy shadowed topologies, First Node Death occurs at round $\approx 110$, validating the realistic stress of spatial occlusion.

\subsection{Belief Routing Under Telemetry Uncertainty}
Table~\ref{tab:uncertainty} reports 180 runs across reporting intervals $R \in [1, 20]$. In the scarce solar regime ($P_{\text{solar}} = 0.25\text{ mJ}$), conservative Belief LCB-Maximin (M4) closes $55.4\%$ of the oracle gap, while Belief DP (M3) closes $81.5\%$ in balanced regimes. The 24-bin histogram reduces Expected Calibration Error by over $40\%$ compared to stale scalar reporting ($0.023$ vs $0.039$).

\subsection{DSU Detour Speedup and Scalability}
Over 500 fault injection trials (\texttt{run\_dsu\_benchmark.py}), DSU local detour repair sustained a 100.0\% packet delivery ratio up to a 30\% relay failure rate, executing in $884.73\ \mu\text{s}$ ($0.88\text{ ms}$) compared to $3443.04\ \mu\text{s}$ ($3.44\text{ ms}$) for full 3D Time-DP recomputation, delivering a $3.9\times - 6.1\times$ latency reduction while avoiding packet drops.

Network scalability benchmarks ($N = 50 \to 500$, \texttt{run\_scalability\_benchmark.py}) confirm theoretical complexity bounds: Dijkstra scales from $0.11\text{ ms} \to 1.14\text{ ms}$, while 3D Time-DP scales from $2.25\text{ ms} \to 7.76\text{ ms}$ (sub-$8\text{ ms}$), confirming real-time embedded feasibility on 32-bit microcontrollers. DSU detour splicing scales from $0.41\text{ ms} \to 6.54\text{ ms}$, preserving its runtime advantage.

\subsection{Real-World Case Study: Great Duck Island Topology}
To assess performance under empirical habitat telemetry parameters, we simulated the Great Duck Island (GDI) deployment ($N=32$ motes, $200\text{m} \times 100\text{m}$ terrain, 60\% spruce canopy shade, 500 rounds; \texttt{run\_real\_world\_case\_study.py}). Conventional LEACH experiences First Node Death at round 433 ($\text{FND}=433$, 31/32 motes alive, $3.4476\text{ J}$ residual), and EH-LEACH at round 445 ($\text{FND}=445$, 31/32 alive, $3.4383\text{ J}$). The proposed Time-Augmented DP + DSU sustains 100\% mote survival with zero node deaths throughout the full 500-round mission ($\text{FND} \ge 500$, 32/32 motes alive), delivering a $+15.5\%$ mission lifetime extension over LEACH while preserving the highest final residual energy ($3.5936\text{ J}$).

\section{Threats to Validity \& Limitations}
\begin{enumerate}
    \item \textbf{Receiver Dissipation ($E_{\text{rx}}$):} Single-hop direct routing consumes less total energy than multi-hop relaying for $d < 87.7\text{ m}$ under clear line-of-sight. Multi-hop routing should be engaged selectively when distance exponents ($d^4$) or severe obstacles dictate.
    \item \textbf{Timescale Decoupling:} We assume transmission delays occur at millisecond resolution while solar variation occurs at minute/hour scales, synchronized by discrete slot offsets.
    \item \textbf{MAC Scheduling:} Assumes contention-free TDMA intra-cluster slots without buffer queue overflow.
\end{enumerate}

\section{Conclusion}
This paper presented a lightweight, provably bounded algorithmic framework for Energy-Harvesting WSNs. By integrating 3D Time-Augmented Dynamic Programming ($dp[v][h][t]$), Disjoint-Set Union (DSU) live detours, 24-bin discrete histogram belief filtering, and realistic first-order radio physics ($E_{\text{rx}}$), the proposed system resolves the classical false rejection of recharging nodes, achieves live fault tolerance in $O(\text{deg}(u)\alpha(V))$ time, and closes up to $81.5\%$ of the oracle delivery gap under stale telemetry. Future research will explore Contiki-NG embedded deployments on physical TelosB motes.

\section*{Acknowledgment}
The authors thank the faculty and laboratory staff of the School of Electrical Engineering (SELECT) at Vellore Institute of Technology (VIT), Vellore, for providing computational infrastructure and technical support during this research project.

\bibliographystyle{IEEEtran}
\begin{thebibliography}{15}

\bibitem{akyildiz2002survey}
I.~F. Akyildiz, W.~Su, Y.~Sankarasubramaniam, and E.~Cayirci, ``Wireless sensor networks: A survey,'' \textit{Computer Networks}, vol.~38, no.~4, pp. 393--422, Mar. 2002.

\bibitem{sah2020renewable}
D.~K. Sah and T.~Amgoth, ``Renewable energy harvesting schemes in wireless sensor networks: A survey,'' \textit{Information Fusion}, vol.~63, pp. 195--213, Nov. 2020.

\bibitem{heinzelman2002application}
W.~R. Heinzelman, A.~Chandrakasan, and H.~Balakrishnan, ``An application-specific protocol architecture for wireless microsensor networks,'' \textit{IEEE Transactions on Wireless Communications}, vol.~1, no.~4, pp. 660--670, Oct. 2002.

\bibitem{sudevalayam2011energy}
S.~Sudevalayam and P.~Kulkarni, ``Energy harvesting sensor nodes: Survey and implications,'' \textit{IEEE Communications Surveys \& Tutorials}, vol.~13, no.~3, pp. 443--461, 3rd Quart. 2011.

\bibitem{sharma2020energy}
S.~K. Sharma, A.~S.~M.~Z. Kausar, and C.~H. Foh, ``Energy harvesting wireless sensor networks: Routing protocols and challenges,'' \textit{IEEE Communications Surveys \& Tutorials}, vol.~22, no.~4, pp. 2988--3013, 4th Quart. 2020.

\bibitem{han2022novel}
B.~Han, J.~Zhao, and L.~Zhou, ``A novel adaptive cluster-based routing protocol for energy-harvesting wireless sensor networks,'' \textit{MDPI Sensors}, vol.~22, no.~4, art. no. 1564, Feb. 2022.

\bibitem{naeem2021dare}
A.~Naeem, M.~A. Javed, and M.~Z. Shakir, ``DARE-SEP: Distance-aware residual energy-efficient SEP for heterogeneous wireless sensor networks,'' \textit{IEEE Transactions on Green Communications and Networking}, vol.~5, no.~2, pp. 1046--1056, Jun. 2021.

\bibitem{sah2023harvested}
D.~K. Sah and T.~Amgoth, ``Harvested energy prediction technique for solar-powered wireless sensor networks,'' \textit{IEEE Sensors Journal}, vol.~23, no.~3, pp. 3208--3218, Feb. 2023.

\bibitem{tan2021throughput}
R.~Tan, G.~Xing, B.~Liu, and J.~Yao, ``Throughput maximization in time-expanded graphs for energy-harvesting wireless sensor networks,'' \textit{ACM Transactions on Sensor Networks}, vol.~17, no.~3, art. no. 29, pp. 1--28, Jun. 2021.

\bibitem{kim2021energy}
J.~Kim and S.~Park, ``Energy-aware opportunistic routing for energy-harvesting wireless sensor networks,'' \textit{IEEE Internet of Things Journal}, vol.~8, no.~7, pp. 5128--5140, Apr. 2021.

\bibitem{aref2023deep}
M.~A. Aref and S.~K. Jayaweera, ``Deep reinforcement learning-based energy-efficient routing in energy harvesting IoT networks,'' \textit{IEEE Transactions on Cognitive Communications and Networking}, vol.~9, no.~3, pp. 689--702, Jun. 2023.

\bibitem{rathee2021qos}
M.~Rathee, S.~Kumar, A.~H. Gandomi, and F.~J. Garcia-Penalvo, ``QoS-aware energy balancing secure routing in wireless sensor networks using ant colony optimization,'' \textit{IEEE Transactions on Engineering Management}, vol.~70, no.~5, pp. 1832--1846, May 2021.

\bibitem{zhang2021energy}
Y.~Zhang, L.~Zhou, and H.~Wang, ``An energy-harvesting-aware clustering routing protocol for wireless sensor networks,'' \textit{IEEE Access}, vol.~9, pp. 30568--30579, Feb. 2021.

\bibitem{sharma2020energy_tax}
S.~K. Sharma, A.~S.~M.~Z. Kausar, and C.~H. Foh, ``Energy harvesting wireless sensor networks: Routing protocols and challenges,'' \textit{IEEE Communications Surveys \& Tutorials}, vol.~22, no.~4, pp. 2988--3013, 2020.

\bibitem{mainwaring2002wireless}
A.~Mainwaring, D.~Culler, J.~Polastre, R.~Szewczyk, and J.~Anderson, ``Wireless sensor networks for habitat monitoring,'' in \textit{Proc. 1st ACM Int. Workshop Wireless Sensor Netw. Appl. (WSNA)}, Atlanta, GA, USA, 2002, pp. 88--97.

\end{thebibliography}

\end{document}
